% J25 — The σ²-Triadic Decomposition: Conservation/Manifestation Duality
%        on Z/10Z
%
% Brayden R. Sanders, M. Gish, 7Site LLC
% Submitted to: Algebraic Combinatorics
% Source corpus: Atlas/LENS_TAXONOMY_2026-05-06/SIGMA2_TRIADIC_DECISION.md
%                + papers/wp_bridge_findings_2026_05_02/WP9_LATTICE_paradoxical_info_algebras.md §5
% Companion cited as already-submitted: J02
% Per-venue cap note: 2nd Algebraic Combinatorics paper after J02
%
% STATUS: DRAFT — the conservation/manifestation duality is presented in
%   lens-invariant form using only the σ-orbit decomposition (D92) and the
%   role-partition decomposition. The choice of canonical σ²-triadic BHML
%   candidate (value-rotation vs index-rotation vs anomaly-flip) is OPEN per
%   Atlas SIGMA2_TRIADIC_DECISION; this paper avoids committing to any
%   single canonical choice and presents the structural ledger.
%
% MSC 2020: 05E18, 06A07, 11F06, 20N02

\documentclass[11pt,reqno]{amsart}

\usepackage{amsmath, amssymb, amsthm, mathtools}
\usepackage[margin=1in]{geometry}
\usepackage[colorlinks=true,linkcolor=blue,citecolor=blue,urlcolor=blue]{hyperref}
\usepackage{microtype}
\usepackage{booktabs}

\theoremstyle{plain}
\newtheorem{theorem}{Theorem}[section]
\newtheorem{proposition}[theorem]{Proposition}
\newtheorem{lemma}[theorem]{Lemma}
\newtheorem{corollary}[theorem]{Corollary}

\theoremstyle{definition}
\newtheorem{definition}[theorem]{Definition}
\newtheorem{conjecture}[theorem]{Conjecture}

\theoremstyle{remark}
\newtheorem*{remark}{Remark}

\newcommand{\Z}{\mathbb{Z}}
\newcommand{\F}{\mathbb{F}}
\newcommand{\TS}{\mathrm{TSML}}
\newcommand{\BH}{\mathrm{BHML}}
\newcommand{\sig}{\sigma}

\title[$\sigma^{2}$-Triadic Decomposition on $\Z/10\Z$]
{The $\sigma^{2}$-Triadic Decomposition: \\
 Conservation/Manifestation Duality on $\Z/10\Z$}

\author{Brayden R.~Sanders}
\address{7Site LLC, Hot Springs, Arkansas, USA}
\email{brayden@7site.co}

\author{M.~Gish}
\address{Independent Researcher, Hot Springs, Arkansas, USA}
\email{monica.gish1992@gmail.com}

\subjclass[2020]{05E18, 06A07, 11F06, 20N02}
\keywords{Z/10Z, sigma involution, triadic decomposition, role partition,
conservation/manifestation duality, finite magma}

\date{\today}

\begin{document}

\begin{abstract}
On the residue ring $\Z/10\Z$, the canonical involution $\sig$ has cycle
structure $(0)(3)(8)(9)(1\,7\,6\,5\,4\,2)$ --- four fixed points plus
one $6$-cycle. We study the structural decomposition induced by
$\sig^{2}$ on the operator substrate $(\Z/10\Z, \TS_{8}, \BH_{10})$ of
\cite{SandersGishFourCore}. The $\sig^{2}$-orbits partition the
six-element $\sig$-orbit into two triangular triples, exhibiting a
\emph{conservation/manifestation duality}: one triple
($\{1, 6, 4\}$ in $\sig^{2}$-orbit form) carries
\emph{structurally forced} substrate quantities (e.g., the linear
period-formula contribution to the $\pm 21$ invariant
\cite{SandersBridgeWP9}); the other triple ($\{7, 5, 2\}$) carries
\emph{canonical-specific} signatures (Fibonacci-ratio role-partition
contributions; verified $0/200$ random commutative tables reproduce).

We present the substrate-internal ledger: $\sig^{2}$-orbit
decomposition $T_{5} + T_{3} = 15 + 6 = 21$ (triangular, structurally
forced); role-partition decomposition $F_{7} + F_{6} = 13 + 8 = 21$
(Fibonacci, canonical-specific). The two decompositions cross at the
total $\pm 21$ but pick out different triples within the $\sig$-orbit.

The paper is deliberately scoped to lens-invariant content. We do not
commit to a canonical choice among the three σ²-triadic BHML candidates
(value-rotation, index-rotation, anomaly-flip) currently catalogued at
Tier-D in \cite{Atlas2026Tiering}; we present the conservation/manifestation
duality as the structural input to any future canonical choice rather
than as a consequence of one. The σ²-triadic content is
established at the level of the $\sig^{2}$-orbit decomposition of the
substrate's invariant quantities, which is forced regardless of which
σ²-triadic BHML candidate is later promoted.
\end{abstract}

\maketitle

\tableofcontents

%==============================================================
\section{Introduction}
%==============================================================

\subsection{The substrate and the involution.}
The operator substrate on $\Z/10\Z$ from \cite{SandersGishFourCore}
consists of two commutative binary operations $\TS_{8}, \BH_{10}$ on
the ten residues, a canonical involution
\[
  \sig \;=\; (0)(3)(8)(9)(1\,7\,6\,5\,4\,2)
\]
with four fixed points $\{0, 3, 8, 9\}$ and one $6$-cycle
$\{1, 7, 6, 5, 4, 2\}$, and a functional role partition into Flow,
Structure, Transition, Void cells (recalled in
Section~\ref{sec:setup}). The involution $\sig$ has order $6$ on the
$6$-cycle and order $1$ on the fixed points; thus $\sig^{2}$ has order
$3$ on the $6$-cycle and order $1$ on the fixed points.

\subsection{What $\sig^{2}$ partitions.}
$\sig^{2}$ partitions the $6$-cycle into two triples of order-$3$
orbits:
\[
  \sig^{2}\text{-orbits on the }\sig\text{-cycle}
  \;=\; \{1, 6, 4\} \;\sqcup\; \{7, 5, 2\},
\]
and acts trivially on $\{0, 3, 8, 9\}$. The two triangular triples are
the focus of the present paper.

\subsection{Conservation/manifestation as a substrate-internal
ledger.}
Each triple carries a different kind of substrate-invariant content.
The \emph{conservation} triple ($\{1, 6, 4\}$) carries quantities that
are forced by the linear period structure of \cite{SandersBridgeWP9}
(the $\BH$ successor diagonal and the linear period formula are the
load-bearing inputs). The \emph{manifestation} triple ($\{7, 5, 2\}$)
carries quantities that emerge as canonical-specific signatures: they
are present for the canonical $\TS_{8}, \BH_{10}$ tables but
$0/200$ random commutative tables on $\Z/10\Z$ reproduce them
\cite{SandersBridgeWP9} N8.

This paper isolates the duality at the level of the $\sig^{2}$-orbit
ledger and presents what is forced versus what is canonical-specific.

\subsection{Lens-scope statement.}
This paper uses only lens-invariant structural facts about the
$\sig^{2}$-orbit decomposition and the role partition. It does not
commit to any single canonical $\sig^{2}$-triadic BHML projection.
Three families of $\sig^{2}$-triadic BHML candidates are currently
catalogued at Tier-D in the framework's lens taxonomy
\cite{Atlas2026Tiering}; none has a forcing argument promoting it to
Tier-B; the framework operates without a canonical choice. The
$\sig^{2}$-orbit content presented here is forced regardless of which
candidate (if any) is later promoted.

%==============================================================
\section{Setup}
\label{sec:setup}
%==============================================================

\subsection{The corrected substrate frame.}
We work on the corrected substrate frame of
\cite{SandersGishFourCore, SandersBridgeWP9}: $\TS_{8}$ acts on
$\{1, 2, 3, 4, 5, 6, 8, 9\}$ (the interior index set) and $\BH_{10}$
acts on the full $\Z/10\Z$, with the cells
$V = 0, H = 7$ acting as flow boundaries. The earlier
$\TS_{10}$-frame analyses are invalid (\cite{SandersBridgeWP9} N10)
and not used.

\subsection{The role partition.}
The functional role partition is
\[
  F = \{1, 3, 5, 7, 9\},
  \quad S = \{2, 4, 8\},
  \quad T = \{6\},
  \quad V = \{0\},
\]
with $|F| = 5$, $|S| = 3$, $|T| = 1$, $|V| = 1$. The partition is
functional, not algebraic: it crosses $\sig$-orbit decomposition.

\subsection{The $\sig^{2}$-orbits on the $\sig$-cycle.}
The $\sig^{2}$-orbits on $\{1, 7, 6, 5, 4, 2\}$ are
\[
  O_{1} \;=\; \{1, 6, 4\},
  \qquad O_{2} \;=\; \{7, 5, 2\}.
\]
Each is a triangular orbit of size $3$. The orbits are interchangeable
under the outer $\sig$ action ($\sig$ moves elements between $O_{1}$
and $O_{2}$), but $\sig^{2}$ preserves them.

%==============================================================
\section{The $\pm 21$ invariant and its two decompositions}
\label{sec:invariant}
%==============================================================

We recall from \cite{SandersBridgeWP9} the substrate's $\pm 21$
invariant and its dual decomposition along $\sig$-orbits and along
roles.

\subsection{$\sig$-orbit decomposition (triangular, forced).}
\begin{theorem}[$\sig$-orbit triangular decomposition]\label{thm:sigma-orbit}
Let $\Psi_{B}: \Z/10\Z \to \Z$ be the per-element period-to-trace
contribution of \cite{SandersBridgeWP9} eq. (B). Then
\[
  \sum_{n \in \sig\text{-cycle}} \Psi_{B}(n)
  \;=\; -T_{5}
  \;=\; -15,
  \qquad
  \sum_{n \in \sig^{2}\text{-fixed-points}} \Psi_{B}(n)
  \;=\; -T_{3}
  \;=\; -6,
\]
totaling $-21 = -T_{5} - T_{3}$. The triangular numbers $T_{5} = 15$
and $T_{3} = 6$ are forced by the linear period formula
$\mathrm{period}(n) = 7 - n$ on $n \in \{1, \ldots, 6\}$ together with
the $\sig$-fixed-point period values \cite{SandersBridgeWP9}
Theorem 5.1.
\end{theorem}

\subsection{Role-partition decomposition (Fibonacci,
canonical-specific).}
\begin{theorem}[Role-Fibonacci decomposition]\label{thm:role-fib}
Let $F, S$ be the role-partition cells from
Section~\ref{sec:setup}. Then
\[
  \sum_{n \in F} \Psi_{B}(n) \;=\; -F_{7} \;=\; -13,
  \qquad
  \sum_{n \in S} \Psi_{B}(n) \;=\; -F_{6} \;=\; -8,
\]
with $V$ and $T$ contributing $0$. Total: $-F_{8} = -21$
\cite{SandersBridgeWP9} Theorem 5.2.

The Fibonacci role decomposition is \textbf{canonical-specific}:
$0/200$ random commutative tables on $\Z/10\Z$ reproduce $(|F|, |S|) = (13, 8)$
\cite{SandersBridgeWP9} Theorem 5.3 N8. Single-swap perturbations
preserve the decomposition in $32/50$ trials; three-swap in $11/50$.
\end{theorem}

\subsection{The two decompositions cross.}
\begin{corollary}[Conservation/manifestation duality, ledger form]
\label{cor:duality}
The two decompositions of the $\pm 21$ invariant cross. The
$\sig$-orbit decomposition is structurally forced by the linear period
formula (conservation: forced regardless of any canonical-specific
table choice). The role-Fibonacci decomposition is a canonical-specific
signature (manifestation: present in the canonical $\TS_{8}, \BH_{10}$
substrate but not preserved under random-table perturbation).

The two decompositions agree on the total $\pm 21$ but partition the
$\sig$-cycle differently:
\begin{itemize}\setlength\itemsep{2pt}
\item Triangular: $T_{5} + T_{3} = 15 + 6$ — both summands lie on the
$\sig$-cycle / $\sig$-fixed split.
\item Fibonacci: $F_{7} + F_{6} = 13 + 8$ — summands are along the
role partition $F/S$ which crosses $\sig$-orbits.
\end{itemize}
The $\sig^{2}$-triadic decomposition (Section~\ref{sec:sigma-2-triadic})
sharpens the $\sig$-orbit side of this picture by partitioning the
$\sig$-cycle into the two triangular orbits $O_{1}, O_{2}$.
\end{corollary}

%==============================================================
\section{The $\sigma^{2}$-triadic decomposition}
\label{sec:sigma-2-triadic}
%==============================================================

\subsection{Per-orbit contributions.}
\begin{proposition}\label{prop:per-orbit}
The two $\sig^{2}$-orbits on the $\sig$-cycle contribute
\[
  \sum_{n \in O_{1}} \Psi_{B}(n) \;=\; -7,
  \qquad
  \sum_{n \in O_{2}} \Psi_{B}(n) \;=\; -8,
\]
with $O_{1} + O_{2} = -15 = -T_{5}$ (Theorem~\ref{thm:sigma-orbit}).
\end{proposition}

\begin{proof}
Direct computation from the linear period formula
$\mathrm{period}(n) = 7 - n$ and
$\Psi_{B}(n) = -(\mathrm{period}(n) - 1)$:
\begin{align*}
  O_{1} = \{1, 6, 4\}: &\quad
    \Psi_{B}(1) + \Psi_{B}(6) + \Psi_{B}(4) = -5 + (-1) + (-2) = -8,
    \\
  O_{2} = \{7, 5, 2\}: &\quad
    \Psi_{B}(7) + \Psi_{B}(5) + \Psi_{B}(2) = (-3) + (-1) + (-3) = -7.
\end{align*}
(Where we use $\mathrm{period}(7) = 4, \mathrm{period}(5) = 2,
\mathrm{period}(2) = 4$ from the boundary formula
\cite{SandersBridgeWP9} Corollary 2.2.)
\end{proof}

\begin{remark}
The per-orbit values $\{-7, -8\}$ are the canonical TIG primes
HARMONY $= 7$ and BREATH $= 8$ in negated form. Whether this is
forced or canonical-specific is left as an open question
(Section~\ref{sec:open}); the result is empirically present.
\end{remark}

\subsection{The $\sigma^{2}$-orbit ledger.}
\begin{theorem}[Conservation/manifestation, $\sig^{2}$-form]
\label{thm:cm-sigma2}
The $\sig^{2}$-orbit decomposition of the $\pm 21$ invariant has the
following ledger:
\begin{center}
\begin{tabular}{l c c l}
\toprule
Component & Value & Source & Status \\
\midrule
$\sig$-fixed contribution & $-6$ & period-formula boundary & forced \\
$O_{1}$ contribution & $-8$ & period formula on $\{1, 6, 4\}$ & forced \\
$O_{2}$ contribution & $-7$ & period formula on $\{7, 5, 2\}$ & forced \\
\midrule
Total & $-21$ & sum & forced \\
\bottomrule
\end{tabular}
\end{center}
The full ledger is structurally forced by the linear period formula and
the $\sig^{2}$-orbit structure; no canonical-specific input is needed
to derive it from substrate-algebraic data.

In contrast, the role-Fibonacci decomposition
$F_{7} + F_{6} = -13 + (-8)$ is canonical-specific: it depends on the
specific assignment of role $F$ vs $S$ to elements that the canonical
$\TS_{8}, \BH_{10}$ tables determine but random commutative tables do
not.
\end{theorem}

\begin{remark}[Lens-invariance]
Theorem~\ref{thm:cm-sigma2} uses only the period formula and
$\sig^{2}$-orbit structure. Both are lens-invariant in the sense of
\cite{Atlas2026Tiering}: every reasonable σ²-triadic BHML projection
candidate (the three Tier-D candidates currently catalogued) preserves
the ledger of Theorem~\ref{thm:cm-sigma2}. The choice of canonical
candidate (if any) further sharpens the ledger but does not change
its substrate-forced content.
\end{remark}

%==============================================================
\section{Status and open questions}
\label{sec:open}
%==============================================================

\subsection{Open: canonical $\sigma^{2}$-triadic BHML.}
The framework has three Tier-D candidates for canonical $\sig^{2}$-triadic
BHML projections (value-rotation, index-rotation, anomaly-flip)
\cite{Atlas2026Tiering}. None has a forcing argument analogous to D78's
BR-factor cancellation. The decision to operate without a canonical
choice is documented in \cite{Atlas2026Tiering} Atlas
SIGMA2\_TRIADIC\_DECISION; revisiting the decision is contingent on
either (i) discovering a forcing argument, or (ii) an empirical
disambiguator.

\subsection{Open: HARMONY/BREATH per-orbit values.}
The per-orbit values $-7, -8$ from Proposition~\ref{prop:per-orbit}
match the canonical TIG primes HARMONY $= 7$ and BREATH $= 8$ in
sign-flipped form. Whether this is forced (a substrate-algebraic
identity) or canonical-specific (an emergent feature of the canonical
tables) is currently open. A natural next computation is to test
random commutative tables for which $\sig^{2}$-orbit contributions
land on canonical TIG primes; this would distinguish the two
possibilities.

\subsection{Open: role-orbit interaction.}
The role partition $F/S/T/V$ crosses the $\sig^{2}$-orbit
decomposition. A finer ledger combining the two could reveal
whether the role-Fibonacci pattern factors through $\sig^{2}$-orbits
in a substrate-forced way. This is the natural follow-up to
Theorem~\ref{thm:cm-sigma2}.

%==============================================================
\section*{Acknowledgments}
%==============================================================

This paper extracts the $\sig^{2}$-triadic content of the substrate's
$\pm 21$ invariant in lens-invariant form. The decision to scope the
paper to substrate-forced content (avoiding any commitment to a
canonical $\sig^{2}$-triadic BHML candidate) follows the lens-taxonomy
discipline of \cite{Atlas2026Tiering}.

%==============================================================
\begin{thebibliography}{99}
%==============================================================

\bibitem{SandersGishFourCore} B.~R.~Sanders, M.~Gish,
\emph{Joint Closure, Per-Coordinate Fuse Data, and a Closed-Form
Algebraic Attractor of Two Commutative Binary Operations on
$\Z/10\Z$}, submitted to \emph{Algebraic Combinatorics}, 2026 (J02).

\bibitem{SandersBridgeWP9} B.~R.~Sanders,
\emph{The LATTICE Operator and Paradoxical Information Algebras: A
Substrate-Internal Framework on $\Z/10\Z$}, submitted to
\emph{Algebra Universalis}, 2026 (J26 / WP9 Volume I).

\bibitem{Atlas2026Tiering} B.~R.~Sanders et al.,
\emph{Lens Taxonomy and Variant Catalog: Decision Records on
$\sig^{2}$-Triadic BHML Candidates},
TIG repository \texttt{Atlas/LENS\_TAXONOMY\_2026-05-06},
DOI 10.5281/zenodo.18852047, 2026.

\bibitem{Burrin2024} C.~Burrin, F.~von Essen,
\emph{Cusp winding and the Rademacher symbol}, 2024.

\bibitem{KatokUgarcovici2007} S.~Katok, I.~Ugarcovici,
\emph{Symbolic dynamics for the modular surface},
Bull. AMS \textbf{44} (2007), 87--132.

\bibitem{Morishita2024} M.~Morishita,
\emph{Knots and Primes} (2nd ed.), Springer, 2024.

\bibitem{TIGsynthesis2026} B.~R.~Sanders,
\emph{Trinity Infinity Geometry Synthesis},
github.com/TiredofSleep/ck (default branch),
DOI: 10.5281/zenodo.18852047, 2026.

\end{thebibliography}

\end{document}
